% !PRINTED 521
\weberpage{521}
In the first case, according to § 105, only the two congruences
\[
A \equiv 1, \quad A \equiv 5\ (\mathrm{mod}\ 8),
\]
in the second case only the two congruences
\[
A \equiv 1, \quad A \equiv 3\ (\mathrm{mod}\ 8)
\]
occur, and again in equally many classes of forms. The complete decomposition of the class equation according to the genera always requires, when $m$ is divisible by 8, the adjunction of the roots of \gesperrt{all} the primes dividing $m$, including 2.

One notices that these considerations contain a new proof of the theorems on the existence of the genera proved by \gesperrt{Gauss} and \gesperrt{Dirichlet}, though only for negative discriminants.

\begin{center}
\S\ 139. \gesperrt{Examples.}
\end{center}

For the actual calculation of the splitting of the class equation the formulae of the preceding section are applicable only to a limited extent, because of the too high degree of the transformation equations. We shall afterwards bring into application, in one case, a method that is at least closely related to them. If the class equation is known, the decomposition can usually be found more easily in a direct way, by making an ansatz of known form with indeterminate coefficients; thus from the formulae § 130, (2), (29), (37) one finds the following partial equations, in which the positive sign of the square root corresponds to the principal genus:
\[
m = 50, \quad f_1\!\left(\sqrt{-50}\right) = \sqrt[4]{2}\, x,
\]
\[
x^3 - x^2 = \frac{1+\sqrt5}{2}(x+1),
\]
\[
m = 26, \quad f_1\!\left(\sqrt{-26}\right)^2 = \sqrt2\, x,
\]
\[
x^3 - x^2 = \frac{3+\sqrt{13}}{2}(x+1),
\]
\[
m = 41, \quad z = \frac{f\!\left(\sqrt{-41}\right)^2}{\sqrt2} + \frac{\sqrt2}{f\!\left(\sqrt{-41}\right)^2},
\]
\[
z^2 - \frac{5+\sqrt{41}}{2} z + \frac{7+\sqrt{41}}{2} = 0.
\]
